\section{Knowledge representation: local formal dependencies}
\label{sec:p19-ckr-annex}

\begin{annexitem}{annex:p19.ckr-metric}{thm:p19-ckr-metric}{The corrected Minkowski formula and its scope}
\begin{proof}
Every summand is nonnegative; it is zero exactly when its coordinate
difference is zero. Thus the distance is nonnegative and vanishes
exactly for identical points. Interchanging the arguments leaves
all absolute differences unchanged, proving symmetry.

For the triangle inequality write \(\|a\|_r=(\sum_i|a_i|^r)^{1/r}\)
for finite coordinate vectors. When \(r=1\), the coordinate
inequalities \(|a_i+b_i|\leq|a_i|+|b_i|\) sum directly to
\(\|a+b\|_1\leq\|a\|_1+\|b\|_1\).

For \(r>1\), put \(s=r/(r-1)\), so \(1/r+1/s=1\).
For nonnegative \(u,v\), Young's inequality
\(uv\leq u^r/r+v^s/s\) follows by maximizing
\(uv-u^r/r\) over \(u\geq0\): its derivative is
\(v-u^{r-1}\), its maximum is at \(u=v^{1/(r-1)}\),
and the maximum value is \(v^s/s\). At \(v=0\) the
same inequality is immediate. Normalize nonzero vectors of
nonnegative coordinates by their respective \(r\)- and
\(s\)-norms and sum Young's inequality. This gives the finite
H\"older inequality
\[
 \sum_i u_iv_i\leq
    \left(\sum_i u_i^r\right)^{1/r}
    \left(\sum_i v_i^s\right)^{1/s}.
\]
If either norm is zero all its coordinates vanish and the
inequality is also immediate, so normalization omitted no case.

Put \(c_i=|a_i+b_i|\). Coordinatewise triangle inequality
and the just-proved bound give
\[
\begin{aligned}
 \sum_i c_i^r
 &\leq\sum_i |a_i|c_i^{r-1}+\sum_i|b_i|c_i^{r-1}\\
 &\leq(\|a\|_r+\|b\|_r)
           \left(\sum_i c_i^{(r-1)s}\right)^{1/s}\\
 &=(\|a\|_r+\|b\|_r)\|a+b\|_r^{r-1}.
\end{aligned}
\]
If \(\|a+b\|_r=0\) the triangle conclusion is immediate.
Otherwise divide by its positive \((r-1)\)-power to obtain
the same conclusion. Taking \(a=x-y\), \(b=y-z\)
proves \(d(x,z)\leq d(x,y)+d(y,z)\).

For the excluded exponent take \(r=1/2\), and points
\((0,0),(1,0),(1,1)\). The two successive distances are
one, whereas the distance between the first and third is
\((1+1)^2=4>2\). Thus the positive-exponent formula does
not always give a metric when the assumption is weakened.
Omitting absolute values is already fatal at \(r=1\):
in one dimension \(d(0,1)=-1\) under the uncorrected formula.
These counterexamples identify the source repair rather than
invent a new distance criterion.

Every embedding into this metric therefore obeys the proved
triangle inequality; a numerical dissimilarity exceeding the
two others' sum cannot be preserved exactly. Qualitative labels
alone need not violate it: lengths one, one and two, for example,
are realized by three collinear points. Assigning the labels
to such lengths requires an interpretation, so labels alone
cannot entail a numerical impossibility. No empirical embedding
or similarity estimate is asserted by this proof.
\end{proof}
\end{annexitem}

\begin{annexitem}{annex:p19.ckr-pairs}{thm:p19-ckr-pairs}{Counting the source's distinct pair judgments}
\begin{proof}
For \(X\geq2\), there are \(X\) choices of a first labelled
item and \(X-1\) different choices of a second. This counts
\(X(X-1)\) ordered pairs. Each unordered pair is counted
exactly twice, once in each order, and neither ordering equals
itself reversed because the items differ. Dividing by two gives
the claimed count. For \(X=0\) and \(X=1\), no pair of
different items exists and the formula also equals zero.
Substitution gives \(15\cdot14/2=105\) and
\(20\cdot19/2=190\).

The assumption of one unordered judgment per pair is the step
permitting division by two. If both directions are elicited,
the ordered count applies instead. Repeated ratings multiply
or otherwise alter the workload, and selecting only some pairs
does not require the complete-pair count. Thus the arithmetic
does not assert a unique experimental design or an effort limit.
\end{proof}
\end{annexitem}

\begin{annexitem}{annex:p19.ckr-decoding}{thm:p19-ckr-decoding}{The task-relative information criterion in both directions}
\begin{proof}
If \(q=g\circ e\) and \(e(s)=e(t)\), then
\(q(s)=g(e(s))=g(e(t))=q(t)\). Thus any correct decoder
requires constancy of the task answer over each encoding fiber.

Conversely suppose this constancy holds. For any actually used
code \(r\in e(\mathcal S)\), at least one scene \(s\)
has \(e(s)=r\). Define \(g(r)=q(s)\). Every such scene
has the same \(q\)-value by the hypothesis, so this specifies
a unique value without ambiguity. It therefore defines a function
on the entire stated decoder domain and gives
\(g(e(s))=q(s)\) for every scene. This proves the converse,
not merely a naming of the required decoder.

In the source's two scenes, the stipulated flat bag has the
same seven feature names for each. The task answer is circle
for the first and square for the second. Since those answers
differ while the codes agree, the necessary condition fails.
No deterministic decoder on that unchanged code can return
both correct answers. An added object/property binding, a
different encoded feature or access to the original scene
changes the available information and does not contradict this
claim. The constructed function is a mathematical decoder;
neither its efficient computation nor its use by a human is
a premise or conclusion of the argument.
\end{proof}
\end{annexitem}
